
\section{Architecture Integration}

\subsection{Model Support}

\begin{table}[H]
\centering
\small
\begin{tabular}{llll}
\toprule
\textbf{Architecture} & \textbf{Target layers} & \textbf{Conv1D?} & \textbf{Nested?} \\
\midrule
DistilBERT-base & \texttt{q\_lin, k\_lin, v\_lin, out\_lin} \& No \& No \\
RoBERTa-base & \texttt{query, key, value, output.dense} \& No \& Yes \\
GPT-2 & \texttt{c\_attn, c\_proj} \& Yes \& No \\
Qwen3.5-0.8B & \texttt{q\_proj, k\_proj, v\_proj, o\_proj} \& No \& No \\
\bottomrule
\end{tabular}
\end{table}

By default, AASVD targets \textbf{attention projections only} 
(\texttt{target\_attention\_only = True}). Set \texttt{target\_attention\_only = False} to 
also compress FFN layers.

\subsection{Layer Targeting Rules}

AASVD protects the following layers from compression:

\begin{table}[H]
\centering
\small
\begin{tabular}{ll}
\toprule
\textbf{Pattern} & \textbf{Reason} \\
\midrule
\texttt{lm\_head} \& Output projection tied to embeddings \\
\texttt{classifier} \& Task-specific head, small \\
\texttt{head} \& Same as above \\
\texttt{output\_layer} \& Task-specific \\
\texttt{score} \& GPT-2 classification head \\
\texttt{pooler} \& BERT pooler, not a projection \\
\bottomrule
\end{tabular}
\end{table}

The default targeting uses \texttt{should\_compress\_linear} for additional checks.

\subsection{Handling Conv1D (GPT-2)}

GPT-2 uses HuggingFace's \texttt{Conv1D}, which stores weights as \texttt{(in, out)} 
instead of \texttt{(out, in)}:

\begin{lstlisting}
# Standard PyTorch Linear
W = raw\_w  # shape (out, in)

# GPT-2 Conv1D
W = raw\_w.t()  # transpose to (out, in)
\end{lstlisting}

AASVD detects Conv1D via:

\begin{lstlisting}
is\_conv1d = (
    "conv1d" in cls\_name 
    or "conv1d" in target\_cls 
    or hasattr(module, "nf") 
    or hasattr(target\_mod, "nf")
)
\end{lstlisting}

The transposed weight is used throughout. The factorized \texttt{SVDLinear} 
receives the same input/output shapes as the original Conv1D, so the 
rest of the model's forward pass works unchanged.

\subsection{Handling LoRA-Wrapped Modules}

When the model has been wrapped with LoRA adapters (\texttt{LoRALinear}), AASVD 
extracts the \texttt{base\_layer}:

\begin{lstlisting}
target\_mod = getattr(module, "base\_layer", module)
\end{lstlisting}

The factorized layer replaces \texttt{module.base\_layer}, preserving the outer 
LoRA parameters. This ensures that any existing LoRA adapters remain 
attached and trainable.

\subsection{Replacing a Layer}

Replacement uses \texttt{replace\_module} from the framework:

\begin{lstlisting}
if hasattr(module, "base\_layer"):
    module.base\_layer = svd\_layer
else:
    replace\_module(self.model, name, svd\_layer)
\end{lstlisting}

After replacement, all parameters in fallback layers are frozen 
(\texttt{requires\_grad = False}) to prevent accidental training.

\subsection{Causality and Streaming}

For decoder models (GPT-2, Qwen), AASVD's calibration is compatible with 
causal forward passes:

\begin{itemize}
\item Forward hooks fire during the standard causal forward pass.
\item No future-token leakage occurs because the model's own attention
  masking is respected.

\end{itemize}
\textbf{Verification.} We check that gradient of output at position $i$ w.r.t. 
input at position $j > i$ is zero, both before and after injection.

\subsection{Injection Algorithm}

\textbf{Algorithm 2:} Inject AASVD into a full model
\begin{lstlisting}
Input:  model M, rank r, threshold $\tau$
        calibration\_texts D
Output: compressed model M'

1:  targets $\leftarrow$ [(name, module) for name, module in M if is\_target(module)]
2:  hooks $\leftarrow$ []
3:  for (name, module) in targets do
4:      hooks.append(register\_hook(module, accumulate\_C\_X(name)))
5:  end for
6:
7:  # Calibration forward pass
8:  tok\_inputs $\leftarrow$ tokenize(D)
9:  M.eval()
10: with no\_grad(): M(tok\_inputs)
11:
12: for h in hooks: h.remove()
13:
14: # Precompute C\_X^{1/2}, C\_X^{-1/2} for each layer
15: cov\_cache $\leftarrow$ {}
16: for (name, module) in targets do
17:     cov\_cache[name] $\leftarrow$ compute\_sqrt\_and\_inv(C\_X[name])
18: end for
19:
20: # Compress each layer
21: for (name, module) in targets do
22:     svd\_layer, err, r\_eff $\leftarrow$ compress\_linear(module, cov\_cache[name])
23:     if svd\_layer is not None and err $\le$ $\tau$:
24:         replace(module, svd\_layer)
25:     else:
26:         freeze\_all(module)
27: end for
28:
29: # Free covariances
30: cov\_cache.clear()
31: torch.cuda.empty\_cache()
32:
33: return M'
\end{lstlisting}

\subsection{The `SVDLinear` Module}

\begin{lstlisting}
class SVDLinear(nn.Module):
    def \_\_init\_\_(self, in\_features, out\_features, rank, bias=True, ...):
        super().\_\_init\_\_()
        self.V = nn.Parameter(torch.empty(rank, in\_features))
        self.U = nn.Parameter(torch.empty(out\_features, rank))
        if bias:
            self.bias = nn.Parameter(torch.empty(out\_features))
        if lora\_r \textgreater{} 0:
            self.lora\_A = nn.Parameter(torch.empty(lora\_r, in\_features))
            self.lora\_B = nn.Parameter(torch.zeros(out\_features, lora\_r))
    
    def forward(self, x):
        latent = F.linear(x, self.V)
        out = F.linear(latent, self.U)
        if self.lora\_rank \textgreater{} 0:
            out += F.linear(F.linear(x, self.lora\_A), self.lora\_B) * self.lora\_scaling
        if self.bias is not None:
            out += self.bias
        return out * self.output\_scale
\end{lstlisting}

\subsection{Interaction with LoRA Recovery}

After AASVD compression, the model is passed to the LoRA recovery pipeline:

\begin{enumerate}
\item \textbf{AASVD} compresses $\to$ factorized \texttt{SVDLinear} layers.
\item \textbf{Framework} attaches LoRA adapters (\texttt{lora\_A}, \texttt{lora\_B}).
\item \textbf{Recovery} trains only LoRA parameters.
\item \textbf{Evaluation} uses the fine-tuned model.

\end{enumerate}
\textbf{Note.} AASVD can also instantiate LoRA adapters \textit{inside} the \texttt{SVDLinear} 
module (via the \texttt{lora\_rank} parameter). This is used when the framework 
expects LoRA to be part of the compression module.

\subsection{Registry}

\begin{lstlisting}
@COMPRESSOR\_REGISTRY.register("aasvd")
@COMPRESSOR\_REGISTRY.register("aa\_svd")
@COMPRESSOR\_REGISTRY.register("aa\_svd\_compressor")
class AASVDCompressor(BaseCompressor):
    ...
\end{lstlisting}

\subsection{Failure Modes and Fallbacks}
